% !TEX program = lualatex
% Transparencias de Fundamentos de las Matemáticas II
% Clase 09: Aritmética modular I

\documentclass[aspectratio=169,11pt]{beamer}

\usepackage{fontspec}
\usepackage[spanish,es-nodecimaldot]{babel}
\usepackage{amsmath,amssymb,mathtools}
\usepackage{booktabs}
\usepackage{csquotes}
\usepackage{xcolor}

\usetheme{Madrid}
\usecolortheme{default}
\setbeamertemplate{navigation symbols}{}
\setbeamertemplate{footline}[frame number]

\definecolor{FdMBlue}{RGB}{20,55,100}
\definecolor{FdMRed}{RGB}{130,30,30}
\definecolor{FdMGreen}{RGB}{25,100,60}

\setbeamercolor{structure}{fg=FdMBlue}
\setbeamercolor{title}{fg=white,bg=FdMBlue}
\setbeamercolor{frametitle}{fg=white,bg=FdMBlue}
\setbeamercolor{block title}{fg=white,bg=FdMBlue}
\setbeamercolor{block body}{bg=blue!3}
\setbeamercolor{alerted text}{fg=FdMRed}

\setmainfont{Latin Modern Roman}
\setsansfont{Latin Modern Sans}
\setmonofont{Latin Modern Mono}

\newcommand{\N}{\mathbb{N}}
\newcommand{\Z}{\mathbb{Z}}
\newcommand{\Q}{\mathbb{Q}}
\newcommand{\R}{\mathbb{R}}
\newcommand{\C}{\mathbb{C}}
\newcommand{\Pow}{\mathcal{P}}
\newcommand{\id}{\operatorname{id}}
\newcommand{\mcd}{\operatorname{mcd}}
\newcommand{\mcm}{\operatorname{mcm}}
\newcommand{\im}{\operatorname{Im}}

\title[FdM II -- Clase 09]{Aritmética modular I}
\subtitle{Congruencias y clases de congruencia}
\author{Antonio Falcó Montesinos}
\institute{Universidad CEU Cardenal Herrera}
\date{Fundamentos de las Matemáticas II}

\begin{document}

\begin{frame}
  \titlepage
\end{frame}

\begin{frame}{Objetivos de la clase}
\begin{itemize}
  \item Definir congruencia módulo \(n\).
  \item Interpretar clases de congruencia.
  \item Construir \(\mathbb{Z}/n\mathbb{Z}\).
\end{itemize}
\end{frame}

\begin{frame}{Mapa de la clase}
\tableofcontents
\end{frame}

\section{Congruencias}

\begin{frame}{Congruencias}
\begin{block}{Definición}
\begin{itemize}
  \item Sean \(a,b\in\mathbb{Z}\) y \(n\geq 1\).
  \item Escribimos \(a\equiv b\pmod n\) si \(n\mid(a-b)\).
  \item La congruencia depende del módulo.
\end{itemize}
\end{block}
\end{frame}

\begin{frame}{Congruencias}
\begin{block}{Ejemplos}
\begin{itemize}
  \item \(17\equiv 2\pmod 5\), porque \(5\mid 15\).
  \item \(-1\equiv 4\pmod 5\), porque \(5\mid -5\).
  \item Los enteros congruentes tienen el mismo resto al dividir por \(n\).
\end{itemize}
\end{block}
\end{frame}

\begin{frame}{Congruencias}
\begin{block}{Relación de equivalencia}
\begin{itemize}
  \item La congruencia módulo \(n\) es reflexiva, simétrica y transitiva.
  \item Por tanto, agrupa \(\mathbb{Z}\) en clases.
  \item Estas clases son las clases de congruencia.
\end{itemize}
\end{block}
\end{frame}

\section{Clases}

\begin{frame}{Clases}
\begin{block}{Clase de \(a\)}
\begin{itemize}
  \item \([a]_n=\{x\in\mathbb{Z}:x\equiv a\pmod n\}\).
  \item Todos los elementos de la clase tienen el mismo resto módulo \(n\).
  \item Distintos representantes pueden nombrar la misma clase.
\end{itemize}
\end{block}
\end{frame}

\begin{frame}{Clases}
\begin{block}{Cociente}
\begin{itemize}
  \item \(\mathbb{Z}/n\mathbb{Z}=\{[0]_n,[1]_n,\ldots,[n-1]_n\}\).
  \item Hay exactamente \(n\) clases.
  \item El representante estándar es el resto entre \(0\) y \(n-1\).
\end{itemize}
\end{block}
\end{frame}

\section{Profundización}

\begin{frame}{Profundización}
\begin{block}{Lectura por restos}
\begin{itemize}
  \item Calcular módulo \(n\) es trabajar con restos.
  \item Las clases permiten sustituir un entero por otro congruente.
  \item La sustitución debe respetar el módulo fijado.
\end{itemize}
\end{block}
\end{frame}

\begin{frame}{Profundización}
\begin{block}{Ejemplo}
\begin{itemize}
  \item En módulo \(4\), \([2]_4=\{\ldots,-6,-2,2,6,10,\ldots\}\).
  \item \([6]_4=[2]_4\).
  \item \([1]_4\neq [2]_4\).
\end{itemize}
\end{block}
\end{frame}

\begin{frame}{Profundización}
\begin{block}{Errores frecuentes}
\begin{itemize}
  \item Cambiar de módulo sin avisar.
  \item Confundir una clase con un único entero.
  \item Olvidar que \([a]_n=[b]_n\) si y solo si \(a\equiv b\pmod n\).
\end{itemize}
\end{block}
\end{frame}


\section{Detalle y práctica}

\begin{frame}{Detalle y práctica}
\begin{block}{Congruencia como igualdad de restos}
\begin{itemize}
  \item \(a\equiv b \pmod n\) significa que \(n\mid(a-b)\).
  \item Equivalentemente, \(a\) y \(b\) tienen el mismo resto al dividir por \(n\).
  \item La formulación por divisibilidad suele ser más útil en demostraciones.
\end{itemize}
\end{block}
\end{frame}

\begin{frame}{Detalle y práctica}
\begin{block}{Actividad breve}
\begin{itemize}
  \item Decidir si \(47\equiv 5 \pmod 7\).
  \item Hallar el representante entre \(0\) y \(8\) de \(-23\) módulo \(9\).
  \item Probar que la congruencia módulo \(n\) es una relación de equivalencia.
\end{itemize}
\end{block}
\end{frame}

\begin{frame}{Cierre}
\begin{itemize}
  \item La congruencia es divisibilidad de diferencias.
  \item Las clases de congruencia agrupan enteros por restos.
  \item \(\mathbb{Z}/n\mathbb{Z}\) contiene clases, no enteros sueltos.
\end{itemize}
\end{frame}

\begin{frame}{Trabajo recomendado}
\begin{itemize}
  \item Releer las secciones correspondientes del manual.
  \item Identificar definiciones, símbolos nuevos y enunciados principales.
  \item Preparar dudas y ejemplos para el seminario asociado.
\end{itemize}
\end{frame}

\end{document}
